\subsection{Permutations, Subsets, and Binomial Coefficients}

\begin{definition} \label{def:permutations-and-sign} \hlt{Permutations and Sign}\\
Let $n\geq1$. Let $\mathfrak{S}_n$ denote the set of permutations of $\{1,\ldots,n\}$.\\
The \hlt{sign} of $\sigma\in\mathfrak{S}_n$ is $\sign(\sigma)=1$ if number of pairs $j<k$ with $\sigma(j)>\sigma(k)$ is even, and $\sign(\sigma)=-1$ otherwise.\\
\end{definition}

\begin{definition} \label{def:factorials-binomial-coefficients} \hlt{Factorials and Binomial Coefficients}
\begin{enumerate}[label=(\roman*)]
\setlength{\itemsep}{0pt}
\item The \hlt{factorial} is defined by $0!\defeq1$ and, for $n\in\N$ with $n\geq1$, $n!\defeq n(n-1)\cdots3\cdot2\cdot1$.
\item For integers $0\leq k\leq n$, the \hlt{binomial coefficient} is $\binom{n}{k}\defeq\frac{n!}{k!(n-k)!}$.
\item The symbol $\binom{n}{k}$ is read \hlt{$n$ choose $k$} and counts the $k$-element subsets of an $n$-element set.
\end{enumerate}
\end{definition}

\begin{example} \label{ex:committee-binomial-coefficient} \hlt{Committees}\\
The number of ways to choose a committee of $3$ people from a class of $20$ people is $\binom{20}{3}=\frac{20!}{3!17!}=\frac{20\cdot19\cdot18}{3\cdot2\cdot1}=1140$.
\end{example}

\begin{proposition} \label{prop:binomial-coefficient-edge-and-symmetry} \hlt{Edge Values and Symmetry}\\
For $n\in\N$ and $k\in\Z$ with $0\leq k\leq n$:
\begin{enumerate}[label=(\roman*)]
\setlength{\itemsep}{0pt}
\item \hlt{Edge values}: $\binom{n}{0}=\binom{n}{n}=1$.
\item \hlt{Symmetry}: $\binom{n}{k}=\binom{n}{n-k}$.
\end{enumerate}
\begin{proof}
Let $n\in\N$ and $k\in\Z$ with $0\leq k\leq n$.
\begin{enumerate}[label=(\roman*)]
\setlength{\itemsep}{0pt}
\item Since $0!\defeq1$, $\binom{n}{0}=\frac{n!}{0!n!}=1$ and $\binom{n}{n}=\frac{n!}{n!0!}=1$.
\item Since $0\leq n-k\leq n$, $\binom{n}{n-k}=\frac{n!}{(n-k)!(n-(n-k))!}=\frac{n!}{(n-k)!k!}=\frac{n!}{k!(n-k)!}=\binom{n}{k}$.
\end{enumerate}
\hlt{Structural Commentary}: Choosing the selected $k$ objects is equivalent to choosing the omitted $n-k$ objects.\\
\hlt{Key Insight}: Complements turn $k$-subsets into $(n-k)$-subsets.
\end{proof}
\end{proposition}
